\scanpage{32}

% L0, upper left (slanted)
% slanted
\margin{$*$ \uline{page \unsure{annexe|arrière}}}

% L0, upper right, detached marks
\margin{$\cdot\!\cdot\quad\cdot\ \cdot$}

% L1–5
$\bigcirc\!\!\!{\scriptstyle\genfrac{}{}{0pt}{}{*}{*}}$ Soit $S$ une catégorie avec $\lambda=(\lambda',\lambda'')$ ens de cônes. On lui associe une catégorie $\bar S$ et un foncteur $S\xrightarrow{\beta}\bar S$ \deleted{tel que} \deleted{\unsure{\(\beta\)|p}} en ajoutant, pour tt \deleted{sommet de} cône \inserted{$c':(x\xrightarrow{p_\alpha}(x_\alpha))$} $\in\lambda'$ \deleted{\unsure{$x\to$}}, un objet $\bar x_c$ à $S$, avec \deleted{une de} \inserted{une flèche $x\to\bar x_c$ et} une famille de flèches $\bar x_c\xrightarrow{q_\alpha}x_\alpha$ factorisant les $p_\alpha$ -- et de même dualmt pour les cônes $c''\in\lambda''$.
% L4–7, left
\margin{\diagram{\begin{tikzcd}[ampersand replacement=\&,row sep=small,column sep=small]
 {}\&{}\&x_1\\
 x\arrow[urr]\arrow[r,"i_c"']\arrow[rr]\arrow[drr]\&\bar x\arrow[ur]\arrow[r]\arrow[dr]\&x_2\\
 {}\&{}\&x_3
\end{tikzcd}}}

% L6–9
On dit que $(\bar S,\bar\lambda)$ est déduit de $(S,\lambda)$ en \uline{dédoublant} les sommets des cônes critiques de $S$ ; $\bar\lambda$ est formé des cônes déduits des cônes $c\in\lambda$ en remplaçant le sommet $x_c$ par $\bar x_c$.

% L9–12
Donc pour tt foncteur $u:S-\unsure{\bar E|E}$, avec $E\in\operatorname{Ob}\Delta$, $\exists$ une \uline{unique} extension en $\bar u:\bar S-E$ qui soit $\bar\lambda$-exact, de plus, $u$ est $\lambda$-exact ssi \deleted{\unsure{tt les|ses}} pour tt $c\in\lambda$, $\bar u(i_c)$ est un isomorphisme.

% L13–14
$(*)$ Définition de foncteurs \uline{$\delta$-préexacts} $D.\xrightarrow{u}F$, $D$ muni d'un ens $\delta=(\delta',\delta'')$ de cônes....
% archive sheet number: 31
